\chapter{Background and cost model}
\label{chap:background}

\section{Notation}

Let $D$ be a document composed of $n$ source files and let $p$ denote
its page count. A build is a function that maps $D$ to a PDF; we write
$T(D)$ for its wall-clock duration. The build is decomposed into a fixed
cost $\alpha$ (format loading, package initialisation), a per-page
typesetting cost $\beta$, and a per-run overhead $\gamma_k$ for the
$k$-th pass.

\begin{definition}[Warm cache]
\label{def:warm}
A build is \emph{warm} if every package, font and image required by $D$
is already materialised on the worker that executes it. Otherwise the
build is \emph{cold}.
\end{definition}

\section{A two-term model}

Under \cref{def:warm} the total time of a $K$-pass build is the sum of
the per-pass costs,
%
\begin{equation}
  \label{eq:model}
  T(p, K) \;=\; \alpha \;+\; \sum_{k=1}^{K}\bigl(\gamma_k + \beta_k\, p\bigr)
  \;\approx\; \alpha + K\gamma + \beta p ,
\end{equation}
%
where the approximation assumes that the per-page cost is stable across
passes. The assumption holds well in practice because the second and
later passes differ from the first only in the resolution of forward
references.

For incremental builds the model changes shape. Only the fraction
$\rho \in [0,1]$ of pages touched since the previous build is
re-typeset, so
%
\begin{align}
  T_{\mathrm{inc}}(p, K) &= \alpha_{\mathrm{inc}} + K\gamma + \beta \rho p ,
  \label{eq:inc} \\[2pt]
  \rho &= \frac{\lvert \mathcal{P}_{\text{dirty}} \rvert}{p} ,
  \qquad
  \mathcal{P}_{\text{dirty}} = \bigl\{\, \pi \in D : h(\pi) \neq h'(\pi) \,\bigr\},
  \label{eq:rho}
\end{align}
%
with $h$ and $h'$ the content hashes before and after the edit.

\begin{theorem}[Crossover]
\label{thm:crossover}
Caching outperforms incremental recompilation only while
$\alpha - \alpha_{\mathrm{inc}} > \beta p\,(1-\rho)$. For fixed $\rho$
the right-hand side grows linearly in $p$, so there exists a page count
$p^{\star}$ beyond which incremental recompilation is always faster.
\end{theorem}

\begin{proof}
Subtract \cref{eq:inc} from \cref{eq:model} and solve the resulting
inequality for $p$. The $K\gamma$ terms cancel.
\end{proof}

\begin{lemma}
\label{lem:bound}
If $\rho \le \tfrac{1}{2}$ then $T_{\mathrm{inc}} \le \tfrac{1}{2}T + \alpha_{\mathrm{inc}}$.
\end{lemma}

\begin{keybox}[What the model predicts]
\Cref{thm:crossover} says the advantage of a warm cache is a constant,
while the advantage of incremental work grows with document size. Long
documents should therefore favour incremental builds even when the cache
is perfectly warm. \Cref{chap:results} tests exactly this.
\end{keybox}

\section{Related work}

Build-time modelling for document systems borrows directly from the
compiler literature, and the canonical description of the multi-pass
structure that makes \cref{eq:model} necessary remains
\autocite{knuth1984texbook,lamport1994latex}. Practical package-level
detail is collected in \autocite{mittelbach2023companion}. Scheduling of
shared build workers, which we treat as out of scope, is addressed by
\textcite{example2026conference}.
